\section{Introduction}\label{sec:introduction}
The reciprocal chiral exchange candidate couples a finite native state, a periodic carrier and a conservative receiver. Its continuous exchange flow has an exact storage account, general projection ceilings and a receiving-and-recovery certificate for a particular coherent empty preparation~\cite{rx,pc,rr}. The native model also has an exact discrete map. Inserting that map raises a separate question: which changes in native storage remain attributable to exchange, and what can be inferred about the interval between native updates?

We keep the already defined native-first hybrid~\cite[Section 6]{rx}. The native map acts once, and the exact exchange flow then acts for a fixed symbolic interval $h>0$. The native coefficients do not scale with $h$. Carrier and receiver states are retained at each jump. Neither a physical clock nor an alternative ordering is supplied.

The principal positive result is Theorem~\ref{thm:initial}: the actual synchronized post-native state has nonzero drive for every phase in the stated launch interval and therefore initially supplies the carrier and receiver. Its short-time interval may depend on phase and the fixed finite coefficients. The exact accounts in Section~\ref{sec:accounts}, the first-jump calculation in Section~\ref{sec:firstjump}, the local ceiling comparison in Section~\ref{sec:capacity} and the limits in Section~\ref{sec:cadence} explain what this initial result does and does not establish.

The global exchange theorem, the native map and hybrid ordering, the exchange invariants, and the restricted retained-receiver ratio are inherited. Their indexed use, the explicit first-jump specialization, the full-phase synchronized source estimate, the local capacity calculation and the qualified limit deductions are the accepted present investigation~\cite{investigation,acceptance}. In particular, failure to satisfy the older launch hypothesis is not a proof that later hybrid receiving or recovery is impossible. It means that a new estimate would be required.

\subsection{Notation and conventions}\label{sec:notation}
Native amplitudes retain the published normalized convention, and angles are in radians. The native coefficients $\epsilon,g,\lambda,k_i$ are distinct from the carrier width $\epsilon_c$, yaw $\delta_g$, readout gain $G$ and operator norm $\|K\|$. The full post-jump vector is $f_{\rm post}$, whereas $f_{{\rm raw},\pm}$ denotes an individual raw port. The cadence is $h$; $h_e$ is a kernel-projection quantity. The carrier length is $L_c$ and the invariant floor is $\mathcal L$.

Relative to the accepted proof, we write $\mathcal L$ for its floor $L$, $\mathbf r_0$ for its zero-phase vector $r$, $\xi$ for its phase-derivative vector $t$, $\zeta$ for its unknown second-order real vector $v_2$, and $q_{\rm cap}$ for its local capacity coefficient $Q$. These changes are one-to-one. Carrier $r_c$, exchange time $t$, prestage $v$ and accumulated work $Q_{A\to B}$ retain their separate meanings. The complete equation and source correspondence accompanies the manuscript.
